% Generated draft: requires mathematical and source review.
\documentclass[11pt]{article}
\usepackage[margin=25mm]{geometry}
\usepackage{fontspec}
\setmainfont{Latin Modern Roman}
\usepackage{amsmath,amssymb,mathtools,unicode-math}
\setmathfont{Latin Modern Math}
\usepackage{enumitem,needspace,xurl,hyperref}
\setlength{\parindent}{0pt}
\setlength{\parskip}{5pt}
\setlength{\emergencystretch}{3em}
\newcommand{\N}{\mathbb{N}}
\newcommand{\Z}{\mathbb{Z}}
\newcommand{\Q}{\mathbb{Q}}
\newcommand{\R}{\mathbb{R}}
\newcommand{\C}{\mathbb{C}}
\newcommand{\F}{\mathbb{F}}
\newcommand{\A}{\mathbb{A}}
\newcommand{\PP}{\mathbb{P}}
\newcommand{\E}{\mathbb{E}}
\newcommand{\eps}{\varepsilon}
\newcommand{\dd}{\,\mathrm{d}}
\begin{document}
\section*{UnsolvedMath ID: 800091}
\textbf{Source problem code:} catalogue entry
\subsection{Definitions and mathematical statement}
Let \(\mathcal F\) be a transversally oriented \(C^0\) foliation of codimension \(q>1\) on a manifold \(M\), and let \(L\) be a compact leaf. Transverse orientation means an orientation of the normal microbundle, not a chosen trivialization. For a basepoint \(x\in L\), the holonomy representation is
\[
\operatorname{hol}_{L,x}\colon \pi_1(L,x)\longrightarrow \operatorname{Germ}_0^+(\R^q),
\]
where \(\operatorname{Germ}_0^+(\R^q)\) is the group of germs at \(0\) of orientation-preserving local homeomorphisms fixing \(0\). Its image is the holonomy group of \(L\).

For an explicit conditional formulation, call \(L\) stable if every open neighborhood \(U\) of \(L\) contains an open saturated neighborhood \(V\) of \(L\):
\[
\forall U\subseteq M\text{ open},\qquad
L\subseteq U\ \Longrightarrow\ \exists V\subseteq M\text{ open and saturated},\quad L\subseteq V\subseteq U.
\]
Here saturated means a union of complete leaves. The supplied source record does not define stability. Consequently, this convention requires confirmation against the intended source meaning; in particular, it does not require the leaves in \(V\) to be compact or impose a local fibration structure.

\textbf{General question.} What conditions on the holonomy of \(L\) guarantee that \(L\) is stable? The source asks for conditions, without specifying a necessary-and-sufficient characterization.

\textbf{Specific question.} Suppose \(q=2\), and write \(\nu_L\) for the oriented normal microbundle of \(L\). Let \(e(\nu_L)\in H^2(L;\Z)\) denote its Euler class. Assume
\[
e(\nu_L)=0,\qquad
\pi_1(L,x)\text{ is amenable},\qquad
H^1(L;\R)=0.
\]
Amenability is understood for the fundamental group as a discrete group, and cohomology is understood as singular cohomology. Does it follow that \(L\) is stable? Equivalently, under the stated stability convention, is the implication
\[
\bigl(e(\nu_L)=0\ \land\ \pi_1(L,x)\text{ amenable}\ \land\ H^1(L;\R)=0\bigr)
\ \Longrightarrow\ L\text{ stable}
\]
valid for every such compact leaf? Vanishing of the Euler class is the stated hypothesis; a trivialization of the normal microbundle is not additionally assumed.

The record supplies no boundary conventions or further hypotheses on the ambient manifold. These omissions, together with the undefined source meaning of stability, require human review before this conditional reconstruction can be treated as a definitive formulation.
\subsection{Short English statement}
The question asks which holonomy conditions ensure stability of a compact leaf in a transversally oriented \(C^0\) foliation of codimension greater than one. Its specific case asks whether, in codimension two, vanishing normal Euler class, amenable fundamental group, and \(H^1(L;\R)=0\) suffice. The draft explicitly uses stability by arbitrarily small saturated neighborhoods; the intended source definition remains to be confirmed.
\subsection{Sources}
\begin{itemize}
\item Supplied local UnsolvedMath registry record, permanent ID \(800091\), exact source problem code \texttt{catalogue entry}, titled \emph{Foliations and Diffeomorphism Groups --- Question 32}. The registry attributes the problem to Nariman and Yazdi, with proposed year 2025, and labels it \texttt{open}. This is a supplied status label, not an independently verified mathematical assessment.
\item Nariman and Yazdi, \emph{Problem List on Foliations and Diffeomorphism Groups} (2025), Question 32. Supplied source URL: \url{https://arxiv.org/abs/2503.18274}. The provenance records access on 2026-10-16 and extraction by \texttt{arxiv\_latex}; these are supplied provenance claims, not a claim of independent source access or verification here.
\item The supplied registry record was created and updated on 2026-10-16 and identifies upstream revision \texttt{d371ee917ac9647fbb06cd6fc529deec6f3dda82}. No additional literature references or independently verified literature assessment are supplied.
\end{itemize}
\end{document}
